\section{Najważniejsze wyniki przeprowadzonych badań}
\label{sec:NajwazniejszeWyniki}

Wyniki przeprowadzonych badań wskazują, że wykorzystanie obliczeń równoległych pozwalało na skrócenie czasu sortowania,
jednak uzyskiwane korzyści różniły się w~zależności od badanej implementacji, liczby jednostek wykonawczych oraz charakterystyki danych wejściowych. 

Dla zbioru \texttt{random\_int} o~rozmiarze 1~000~000 elementów najlepsze wyniki czasowe uzyskano dla Bucket Sort, 
dla którego średni czas wykonania wyniósł około 1,20~s przy 16 jednostkach wykonawczych.
W~przypadku Quick Sort najlepszy wynik wyniósł około 1,80~s przy 2 jednostkach, 
dla Merge Sort 2,15~s przy 4 jednostkach, a~dla Sample Sort 2,37~s przy 8 jednostkach.

Najwyższe przyspieszenie odnotowano dla Sample Sort, dla którego wartość ta wyniosła około 2,60. 
Bucket Sort osiągnął przyspieszenie na poziomie około 1,80, natomiast dla Merge Sort i~Quick Sort wyniosło ono odpowiednio 1,59 i~1,21. 
Jednocześnie zwiększanie liczby jednostek wykonawczych nie przekładało się na proporcjonalny wzrost wydajności. 
Przy 32 jednostkach efektywność wyniosła około 0,06 dla Sample Sort, 0,05 dla Bucket Sort i~Merge Sort oraz 0,03 dla Quick Sort.

Dodatkowa diagnostyka wykazała, że liczba zadeklarowanych jednostek wykonawczych 
nie zawsze przekładała się bezpośrednio na liczbę faktycznie tworzonych procesów.
W~przypadku Quick Sort i~Merge Sort liczba procesów zależała również od sposobu podziału danych, 
maksymalnej głębokości zrównoleglenia oraz przyjętych progów uruchamiania wykonania równoległego. 
Dla Quick Sort utworzono odpowiednio 2, 6, 4, 4 oraz 4 procesy przy wykorzystaniu 2, 4, 8, 16 i~32 jednostek wykonawczych, 
natomiast dla Merge Sort liczby te wynosiły odpowiednio 2, 6, 2, 2 oraz 2.

Odmienną charakterystykę zaobserwowano dla Bucket Sort, gdzie liczba utworzonych procesów odpowiadała liczbie zadeklarowanych jednostek wykonawczych.
Sample Sort tworzył natomiast procesy w~dwóch odrębnych fazach przetwarzania, obejmujących sortowanie lokalne oraz sortowanie kubełków.
W~każdej z~faz tworzono liczbę procesów odpowiadającą zadeklarowanej liczbie jednostek, dlatego łączna liczba utworzonych procesów była dwukrotnie większa.

Wyniki pokazują zatem, że zależność pomiędzy liczbą zadeklarowanych jednostek wykonawczych 
a~liczbą tworzonych procesów była uzależniona od sposobu organizacji równoległości w~konkretnym algorytmie.

Przy 8 jednostkach najwyższe średnie obciążenie procesora uzyskano dla Sample Sort około 426\%, a~najniższe dla Quick Sort około 174\%. 
Najmniejsze średnie zużycie pamięci odnotowano dla Merge Sort około 163~MB, natomiast największe dla Sample Sort około 258~MB. 
Przy 32 jednostkach maksymalne zużycie pamięci dla Bucket Sort i~Sample Sort przekraczało 1~GB.

Przeprowadzone eksperymenty pokazały, że charakterystyka danych wejściowych wpływała na badane algorytmy w różnym stopniu. 
Największe różnice w czasie wykonania pomiędzy analizowanymi rodzajami danych odnotowano dla Quick Sort, 
podczas gdy dla pozostałych algorytmów zmienność wyników była na ogół mniejsza.

Zestawienie wyników dla danych całkowitych i~zmiennoprzecinkowych nie wykazało 
jednoznacznej przewagi jednego z tych typów dla wszystkich analizowanych algorytmów. 
Także obecność duplikatów nie wpływała w jednakowy sposób na czas wykonania poszczególnych algorytmów. 
Pokazuje to, że znaczenie charakterystyki danych wejściowych różniło się przede wszystkim w zależności od zastosowanej implementacji.